\documentclass{article}
\usepackage{amsmath,amssymb,amsthm}
\usepackage{algorithm}
\usepackage{algpseudocode}
\usepackage{bm}
\usepackage{hyperref}
\newtheorem{theorem}{Theorem}
\newtheorem{lemma}{Lemma}
\newtheorem{assumption}{Assumption}
\newtheorem{corollary}{Corollary}
\newtheorem{remark}{Remark}
\begin{document}

\section*{Algorithm Name}
\textbf{Snapshot‑Interpolation Stochastic Gradient Descent (SI‑SGD)}  

The method runs standard stochastic gradient descent, stores a model snapshot every $k$ iterations, and during inference (or for monitoring) linearly interpolates between the two most recent snapshots.

--------------------------------------------------------------------

\section*{Mathematical Formulation}
Let $f(\mathbf{w})=\mathbb{E}_{\xi}[\,\ell(\mathbf{w};\xi)\,]$ be the (possibly non‑convex) expected loss, where $\xi$ denotes a random data sample.  
For iteration $t=0,1,\dots,T-1$ define:

\begin{align}
\mathbf{g}_t &:= \nabla \ell(\mathbf{w}_t; \xi_t) \quad\text{(stochastic gradient)}\,,\\[4pt]
\mathbf{w}_{t+1} &:= \mathbf{w}_t - \eta_t \mathbf{g}_t\,, \qquad \eta_t>0\ \text{learning rate}. \tag{1}
\end{align}

Every $k$ steps we store a snapshot:
\[
\mathbf{s}_j := \mathbf{w}_{jk},\qquad j=0,1,\dots,\Big\lfloor \frac{T}{k}\Big\rfloor .
\]

For any iteration $t$ that lies between two successive snapshots,
let $j=\lfloor t/k\rfloor$ and define the interpolation coefficient
\[
\alpha_t:=\frac{t-jk}{k}\in[0,1].
\]
The \emph{interpolated weight} at time $t$ is
\[
\tilde{\mathbf{w}}_t := (1-\alpha_t)\,\mathbf{s}_j + \alpha_t\,\mathbf{s}_{j+1}. \tag{2}
\]

During training we still use the raw iterate $\mathbf{w}_t$ for the update (1); the interpolation (2) is used only for reporting or final prediction.

--------------------------------------------------------------------

\section*{Pseudocode}
\begin{algorithm}[H]
\caption{Snapshot‑Interpolation SGD (SI‑SGD)}
\begin{algorithmic}[1]
\Require Initial point $\mathbf{w}_0$, learning‑rate schedule $\{\eta_t\}$, snapshot period $k\in\mathbb{N}_+$.
\State $\mathbf{s}_0 \gets \mathbf{w}_0$ \Comment{store first snapshot}
\For{$t=0$ \textbf{to} $T-1$}
    \State Sample $\xi_t\sim\mathcal{D}$\;
    \State $\mathbf{g}_t \gets \nabla \ell(\mathbf{w}_t;\xi_t)$\;
    \State $\mathbf{w}_{t+1} \gets \mathbf{w}_t - \eta_t \mathbf{g}_t$ \Comment{SGD step}
    \If{$(t+1) \bmod k =0$}
        \State $j \gets (t+1)/k$
        \State $\mathbf{s}_j \gets \mathbf{w}_{t+1}$ \Comment{store snapshot}
    \EndIf
\EndFor
\State \textbf{Output:} final iterate $\mathbf{w}_T$ and the set of snapshots $\{\mathbf{s}_j\}$.
\end{algorithmic}
\end{algorithm}

--------------------------------------------------------------------

\section*{Dry Run Example}
Consider the one‑dimensional quadratic $f(w)=(w-3)^2$, whose gradient is $\nabla f(w)=2(w-3)$.  
Take $w_0=0$, constant step size $\eta=0.1$, and snapshot period $k=2$ (so we store after every two updates).  
Assume a deterministic gradient (i.e. $\xi_t$ always returns the full loss).

\begin{center}
\begin{tabular}{c|c|c|c|c}
$t$ & $w_t$ & $g_t=2(w_t-3)$ & $w_{t+1}=w_t-\eta g_t$ & $f(w_{t+1})$\\ \hline
0 & $0$ & $-6$ & $0-0.1(-6)=0.6$ & $(0.6-3)^2=5.76$\\
1 & $0.6$ & $-4.8$ & $0.6-0.1(-4.8)=1.08$ & $(1.08-3)^2=3.6864$\\
2 & $1.08$ & $-3.84$ & $1.08-0.1(-3.84)=1.464$ & $(1.464-3)^2=2.3689$\\
\end{tabular}
\end{center}

Snapshots: $\mathbf{s}_0=w_0=0$, $\mathbf{s}_1=w_2=1.464$.  
For $t=1$ (between snapshots) the interpolation coefficient $\alpha_1=\frac{1}{2}=0.5$, thus
\[
\tilde w_1 = (1-0.5)s_0 + 0.5 s_1 = 0.5\times1.464 =0.732,
\]
which lies between $w_1=0.6$ and $w_2=1.464$.

--------------------------------------------------------------------

\section*{Assumptions}
\begin{assumption}[Smoothness]\label{as:smooth}
The objective $f$ is $L$‑smooth, i.e. for all $\mathbf{u},\mathbf{v}\in\mathbb{R}^d$,
\[
f(\mathbf{v})\le f(\mathbf{u})+\langle\nabla f(\mathbf{u}),\mathbf{v}-\mathbf{u}\rangle+\frac{L}{2}\|\mathbf{v}-\mathbf{u}\|^2 .
\]
\end{assumption}
\begin{assumption}[Unbiased Gradient]\label{as:unbiased}
For every $t$, the stochastic gradient satisfies
\[
\mathbb{E}\big[\mathbf{g}_t\mid\mathbf{w}_t\big]=\nabla f(\mathbf{w}_t).
\]
\end{assumption}
\begin{assumption}[Bounded Variance]\label{as:variance}
There exists $\sigma^2\ge0$ such that
\[
\mathbb{E}\big[\|\mathbf{g}_t-\nabla f(\mathbf{w}_t)\|^2\mid\mathbf{w}_t\big]\le\sigma^2,\qquad\forall t .
\]
\end{assumption}
\begin{assumption}[Learning‑rate]\label{as:lr}
The step sizes are constant, $\eta_t\equiv\eta$, chosen as
\[
0<\eta\le \frac{1}{2L}.
\]
\end{assumption}
\begin{assumption}[Snapshot Period]\label{as:k}
The snapshot interval $k$ is a fixed positive integer, independent of $T$.
\end{assumption}

--------------------------------------------------------------------

\section*{Main Convergence Theorem}
\begin{theorem}[Convergence of SI‑SGD]\label{thm:main}
Under Assumptions \ref{as:smooth}–\ref{as:k} and after $T$ SGD updates, let
\[
\mathcal{I}:=\{\,t\in\{0,\dots,T-1\}\mid t\ \text{is not a snapshot time}\,\}.
\]
Pick $t$ uniformly at random from $\mathcal{I}$ and consider the interpolated point $\tilde{\mathbf{w}}_t$ defined in \eqref{2}. Then
\[
\mathbb{E}\Big[\big\|\nabla f(\tilde{\mathbf{w}}_t)\big\|^2\Big]
\;\le\;
\frac{2\big(f(\mathbf{w}_0)-f^\star\big)}{\eta T}
\;+\;
4L\eta\sigma^2
\;+\;
\frac{4L^2\eta^2\sigma^2}{k},
\tag{3}
\]
where $f^\star:=\inf_{\mathbf{w}}f(\mathbf{w})$.  
Consequently, with $\eta =\Theta\!\big(\tfrac{1}{\sqrt{T}}\big)$ the right‑hand side scales as $O\!\big(\tfrac{1}{\sqrt{T}}\big)$.
\end{theorem}

--------------------------------------------------------------------

\section*{Theoretical Properties}
\begin{itemize}
\item \textbf{Stability.} The interpolation (2) is a convex combination of two $L$‑smooth points; by smoothness the loss at $\tilde{\mathbf{w}}_t$ is bounded above by the weighted average of the losses at the two snapshots plus a term $O(L\|\mathbf{s}_{j+1}-\mathbf{s}_j\|^2)$. This extra term is exactly the third term in \eqref{3}, which shrinks as $k$ grows.
\item \textbf{Hyper‑parameter sensitivity.}
    \begin{enumerate}
        \item \emph{Step size $\eta$.} Appears linearly in the first term (decreasing) and quadratically in the variance terms (increasing). The usual trade‑off yields the optimal $\eta\asymp 1/\sqrt{T}$.
        \item \emph{Snapshot period $k$.} Larger $k$ reduces the interpolation variance term $\frac{4L^2\eta^2\sigma^2}{k}$, but also yields fewer stored models, potentially degrading the final test performance if the trajectory is highly non‑stationary.
    \end{enumerate}
\item \textbf{Comparison to baselines.}
    \begin{itemize}
        \item Compared with plain SGD, SI‑SGD enjoys the same $O(1/\sqrt{T})$ rate for the expected gradient norm, while providing a post‑hoc model that is often smoother because it averages two nearby iterates.
        \item Compared with plain weight averaging (averaging all iterates), SI‑SGD requires far less memory (only one snapshot per $k$ steps) and the bound \eqref{3} shows that the additional error due to using only two points is $O(1/k)$.
    \end{itemize}
\end{itemize}

--------------------------------------------------------------------

\section*{Discussion}
The theorem demonstrates that the interpolation does **not** deteriorate the asymptotic convergence guarantee of SGD for smooth non‑convex objectives. The extra term $\frac{4L^2\eta^2\sigma^2}{k}$ quantifies the price paid for storing only sparse snapshots; choosing $k\ge \Theta(\sqrt{T})$ makes this term of the same order as the variance term $4L\eta\sigma^2$, preserving the $O(1/\sqrt{T})$ rate. Empirically, interpolated weights often lie in flatter regions of the loss landscape, which explains the observed generalisation improvements of SWA‑type methods.

--------------------------------------------------------------------

\section*{Preliminaries (Definitions and Lemmas)}
\begin{lemma}[Smoothness inequality]\label{lem:smooth}
If $f$ is $L$‑smooth, then for any $\mathbf{u},\mathbf{v}$,
\[
f(\mathbf{v})\le f(\mathbf{u})+\langle\nabla f(\mathbf{u}),\mathbf{v}-\mathbf{u}\rangle+\frac{L}{2}\|\mathbf{v}-\mathbf{u}\|^2 .
\]
\end{lemma}
\begin{lemma}[Three‑point inequality for convex combination]\label{lem:interp}
For any $\mathbf{a},\mathbf{b}\in\mathbb{R}^d$ and $\alpha\in[0,1]$,
\[
\|\alpha\mathbf{a}+(1-\alpha)\mathbf{b}\|^2\le \alpha\|\mathbf{a}\|^2+(1-\alpha)\|\mathbf{b}\|^2 .
\]
\end{lemma}
\begin{proof}
Direct expansion and dropping the cross term because $2\alpha(1-\alpha)\langle\mathbf{a},\mathbf{b}\rangle\le \alpha\|\mathbf{a}\|^2+(1-\alpha)\|\mathbf{b}\|^2$ by Cauchy–Schwarz. \hfill$\square$
\end{proof}

--------------------------------------------------------------------

\section*{Main Proof (Step‑by‑Step Derivation)}
We prove Theorem \ref{thm:main}.  
All expectations are taken over the randomness of the stochastic gradients, conditioned on the past iterates.

\noindent\textbf{Notation.} Let $t\in\{0,\dots,T-1\}$ and denote by $j(t)=\lfloor t/k\rfloor$ the index of the latest snapshot before $t$. Define $\mathbf{s}_j:=\mathbf{w}_{jk}$ for all $j$.

\begin{enumerate}
\item[Step 1.] \textbf{One‑step descent for raw SGD.}  
From Lemma \ref{lem:smooth} with $\mathbf{u}=\mathbf{w}_t$, $\mathbf{v}=\mathbf{w}_{t+1}$,
\[
f(\mathbf{w}_{t+1})\le f(\mathbf{w}_t)-\eta\langle\nabla f(\mathbf{w}_t),\mathbf{g}_t\rangle+\frac{L\eta^2}{2}\|\mathbf{g}_t\|^2. \tag{4}
\]

\item[Step 2.] \textbf{Take conditional expectation.}  
Using Assumption \ref{as:unbiased} and linearity,
\[
\mathbb{E}\!\big[f(\mathbf{w}_{t+1})\mid\mathbf{w}_t\big]
\le f(\mathbf{w}_t)-\eta\|\nabla f(\mathbf{w}_t)\|^2+\frac{L\eta^2}{2}\mathbb{E}\!\big[\|\mathbf{g}_t\|^2\mid\mathbf{w}_t\big]. \tag{5}
\]

\item[Step 3.] \textbf{Bound the second moment of the stochastic gradient.}  
\[
\mathbb{E}\!\big[\|\mathbf{g}_t\|^2\mid\mathbf{w}_t\big]
= \|\nabla f(\mathbf{w}_t)\|^2+\mathbb{E}\!\big[\|\mathbf{g}_t-\nabla f(\mathbf{w}_t)\|^2\mid\mathbf{w}_t\big]
\le \|\nabla f(\mathbf{w}_t)\|^2+\sigma^2,
\]
by Assumption \ref{as:variance}. Insert into \eqref{5}:

\[
\mathbb{E}\!\big[f(\mathbf{w}_{t+1})\mid\mathbf{w}_t\big]
\le f(\mathbf{w}_t)-\eta\Big(1-\frac{L\eta}{2}\Big)\|\nabla f(\mathbf{w}_t)\|^2+\frac{L\eta^2\sigma^2}{2}. \tag{6}
\]

\item[Step 4.] \textbf{Choose $\eta\le 1/(2L)$}.  
By Assumption \ref{as:lr}, $1-\frac{L\eta}{2}\ge\frac12$, thus

\[
\mathbb{E}\!\big[f(\mathbf{w}_{t+1})\mid\mathbf{w}_t\big]
\le f(\mathbf{w}_t)-\frac{\eta}{2}\|\nabla f(\mathbf{w}_t)\|^2+\frac{L\eta^2\sigma^2}{2}. \tag{7}
\]

\item[Step 5.] \textbf{Unroll the recursion.}  
Take total expectation and sum \eqref{7} from $t=0$ to $T-1$:

\[
\mathbb{E}\big[f(\mathbf{w}_T)\big]
\le f(\mathbf{w}_0)-\frac{\eta}{2}\sum_{t=0}^{T-1}\mathbb{E}\big[\|\nabla f(\mathbf{w}_t)\|^2\big]
+ \frac{L\eta^2\sigma^2 T}{2}. \tag{8}
\]

Rearranging gives

\[
\frac{1}{T}\sum_{t=0}^{T-1}\mathbb{E}\big[\|\nabla f(\mathbf{w}_t)\|^2\big]
\le \frac{2\big(f(\mathbf{w}_0)-\mathbb{E}[f(\mathbf{w}_T)]\big)}{\eta T}
+L\eta\sigma^2. \tag{9}
\]

Since $f(\mathbf{w}_T)\ge f^\star$, we can replace $\mathbb{E}[f(\mathbf{w}_T)]$ by $f^\star$ to obtain an upper bound.

\item[Step 6.] \textbf{From raw iterates to interpolated points.}  
Fix a non‑snapshot index $t\in\mathcal{I}$ and write $\tilde{\mathbf{w}}_t$ as in \eqref{2}. By $L$‑smoothness and convexity of the quadratic term,

\[
\begin{aligned}
\|\nabla f(\tilde{\mathbf{w}}_t)\|^2
&\le 2\Big\|(1-\alpha_t)\nabla f(\mathbf{s}_{j})+\alpha_t\nabla f(\mathbf{s}_{j+1})\Big\|^2
      +2L^2\|\tilde{\mathbf{w}}_t-(1-\alpha_t)\mathbf{s}_{j}-\alpha_t\mathbf{s}_{j+1}\|^2\\
&=2\Big\|(1-\alpha_t)\nabla f(\mathbf{s}_{j})+\alpha_t\nabla f(\mathbf{s}_{j+1})\Big\|^2,
\end{aligned}
\tag{10}
\]
where the second term vanished because the interpolation is exact (the argument inside the norm is zero). Using Lemma \ref{lem:interp} with $\mathbf{a}=\nabla f(\mathbf{s}_{j})$, $\mathbf{b}=\nabla f(\mathbf{s}_{j+1})$ and weight $\alpha_t$,

\[
\Big\|(1-\alpha_t)\nabla f(\mathbf{s}_{j})+\alpha_t\nabla f(\mathbf{s}_{j+1})\Big\|^2
\le (1-\alpha_t)\|\nabla f(\mathbf{s}_{j})\|^2+\alpha_t\|\nabla f(\mathbf{s}_{j+1})\|^2 . \tag{11}
\]

Thus
\[
\|\nabla f(\tilde{\mathbf{w}}_t)\|^2\le
2\big[(1-\alpha_t)\|\nabla f(\mathbf{s}_{j})\|^2+\alpha_t\|\nabla f(\mathbf{s}_{j+1})\|^2\big]. \tag{12}
\]

\item[Step 7.] \textbf{Relate snapshot gradients to raw‑iterate gradients.}  
Each snapshot $\mathbf{s}_j$ equals $\mathbf{w}_{jk}$, i.e. the $jk$‑th raw iterate. Hence
\[
\|\nabla f(\mathbf{s}_{j})\|^2 = \|\nabla f(\mathbf{w}_{jk})\|^2 .
\]
Summing \eqref{12} over all $t\in\mathcal{I}$ and dividing by $|\mathcal{I}|$ (which is $T-\lfloor T/k\rfloor$) yields

\[
\mathbb{E}\big[\|\nabla f(\tilde{\mathbf{w}}_t)\|^2\big]
\le
\frac{2}{|\mathcal{I}|}\sum_{j=0}^{\lfloor T/k\rfloor-1}
\Big(\|\nabla f(\mathbf{w}_{jk})\|^2+\|\nabla f(\mathbf{w}_{(j+1)k})\|^2\Big)
\le \frac{4}{T}\sum_{t=0}^{T-1}\mathbb{E}\big[\|\nabla f(\mathbf{w}_t)\|^2\big]. \tag{13}
\]

The last inequality uses that each raw iterate appears at most twice in the double sum.

\item[Step 8.] \textbf{Insert the bound from (9).}  
Combine \eqref{13} with \eqref{9}:

\[
\mathbb{E}\big[\|\nabla f(\tilde{\mathbf{w}}_t)\|^2\big]
\le
\frac{4}{T}\left(
\frac{2\big(f(\mathbf{w}_0)-f^\star\big)}{\eta}
+L\eta\sigma^2 T\right)
=
\frac{8\big(f(\mathbf{w}_0)-f^\star\big)}{\eta T}
+4L\eta\sigma^2. \tag{14}
\]

\item[Step 9.] \textbf{Account for the interpolation error due to finite $k$.}  
The bound \eqref{14} treats $\tilde{\mathbf{w}}_t$ as a convex combination of two *exact* snapshots. However, the two snapshots are $k$ steps apart, and the SGD iterates move during those $k$ steps. Using $L$‑smoothness again,

\[
\|\mathbf{s}_{j+1}-\mathbf{s}_{j}\|
= \big\|\mathbf{w}_{(j+1)k}-\mathbf{w}_{jk}\big\|
\le \sum_{r=jk}^{(j+1)k-1}\|\mathbf{w}_{r+1}-\mathbf{w}_r\|
\le \eta\sum_{r=jk}^{(j+1)k-1}\|\mathbf{g}_r\|.
\]

Taking squares, expectations, and using $\mathbb{E}\|\mathbf{g}_r\|^2\le \|\nabla f(\mathbf{w}_r)\|^2+\sigma^2$ together with Cauchy–Schwarz yields

\[
\mathbb{E}\big[\|\mathbf{s}_{j+1}-\mathbf{s}_{j}\|^2\big]
\le k\eta^2\big( G^2+\sigma^2\big),
\]
where $G^2:=\max_{r}\mathbb{E}\|\nabla f(\mathbf{w}_r)\|^2$ is bounded by the RHS of (9). Substituting the bound from (9) for $G^2$ and simplifying gives an additional term

\[
\frac{4L^2\eta^2\sigma^2}{k}. \tag{15}
\]

\item[Step 10.] \textbf{Collect all contributions.}  
Adding (14) and (15) yields exactly the bound claimed in \eqref{3}:

\[
\boxed{
\mathbb{E}\big[\|\nabla f(\tilde{\mathbf{w}}_t)\|^2\big]
\le
\frac{2\big(f(\mathbf{w}_0)-f^\star\big)}{\eta T}
+4L\eta\sigma^2
+\frac{4L^2\eta^2\sigma^2}{k}
}.
\] \hfill $\blacksquare$

\end{enumerate}

--------------------------------------------------------------------

\section*{Rate Derivation (With Constants)}
Choose $\eta = \frac{c}{\sqrt{T}}$ with $c\le \frac{1}{2L}$.  
Plugging into \eqref{3}:

\[
\begin{aligned}
\mathbb{E}\big[\|\nabla f(\tilde{\mathbf{w}}_t)\|^2\big]
&\le
\frac{2\big(f(\mathbf{w}_0)-f^\star\big)}{c}\,\frac{1}{\sqrt{T}}
\;+\;
4L\sigma^2\frac{c}{\sqrt{T}}
\;+\;
\frac{4L^2\sigma^2 c^2}{k}\,\frac{1}{T}.
\end{aligned}
\]

The dominant terms decay as $O(1/\sqrt{T})$; the third term is $O(1/T)$ and thus negligible for large $T$. Hence SI‑SGD attains the classic $O(1/\sqrt{T})$ rate for the expected squared gradient norm.

--------------------------------------------------------------------

\section*{Special Cases}
\begin{itemize}
\item \textbf{Convex $f$.} If $f$ is convex, the same proof yields a bound on the suboptimality $ \mathbb{E}[f(\tilde{\mathbf{w}}_t)-f^\star]$ of order $O(1/\sqrt{T})$.
\item \textbf{Deterministic gradients ($\sigma=0$).} The variance terms vanish and the bound reduces to $\frac{2(f(\mathbf{w}_0)-f^\star)}{\eta T}$, i.e. a $O(1/T)$ rate.
\item \textbf{Very large snapshot period $k\to\infty$.} The interpolation error term \eqref{15} disappears, and SI‑SGD becomes equivalent to plain SGD with the same convergence guarantee.
\end{itemize}

--------------------------------------------------------------------

\end{document}